% ==============================================================================
% ACADEMIC RESEARCH POSTER TEMPLATE (tikzposter 3-Column Landscape)
% Author Template by: Nazmul Hasan
% GitHub Repository: https://github.com/Nazmul142965/Attention-In-Poster
% Full Example Demo: https://github.com/Nazmul142965/Attention-In-Poster/tree/main/Assets%20And%20Codes
% ==============================================================================

\documentclass[25pt, a0paper, landscape]{tikzposter}

% ==========================================
% Packages & Font Setup
% ==========================================
\usepackage{iftex}

\ifPDFTeX
  % Default compiler (pdfLaTeX) - works out-of-the-box on Overleaf
  \usepackage[utf8]{inputenc}
  \usepackage[T1]{fontenc}
\else
  % For XeLaTeX or LuaLaTeX (supports custom fonts and non-Latin scripts like Bengali)
  \usepackage{fontspec}
  \setmainfont{FreeSerif}
  % Optional: Bengali script font support
  % \newfontfamily\bangla{FreeSerif}[Script=Bengali]
\fi

\usepackage{amsmath, amssymb}
\usepackage{graphicx}
\usepackage{booktabs}

% ==========================================
% Color & Theme Setup
% ==========================================
% Pre-built themes in tikzposter:
% Default, Rays, Basic, Simple, Envelope, Wave, Board, Autumn, Desert
\usetheme{Board} 
\colorlet{backgroundcolor}{white}

% ==========================================
% Title Information
% ==========================================
\title{Your Research / Project Title Here}
\author{Author Name (e.g., ID: 202XXXXXXX)}
\institute{Department of ..., University / Institution Name}

\begin{document}

\maketitle

% ==========================================
% 3-Column Poster Layout
% ==========================================
\begin{columns}

    % ================= COLUMN 1 =================
    \column{0.33}
    
    \block{1. Overview / Introduction}{
        Write your project introduction, background, and motivation here. 
        Explain the core problem your work solves.
        
        \vspace{0.4cm}
        \begin{itemize}
            \item \textbf{Key Point 1:} Motivation or challenge statement.
            \item \textbf{Key Point 2:} Core approach or solution.
            \item \textbf{Key Point 3:} Summary of contributions.
        \end{itemize}
    }
    
    \block{2. Mathematical Formulation / Methodology}{
        Describe the mathematical foundation and key formulas here.
        
        \vspace{0.4cm}
        \[
        y = f(x; \theta) + \epsilon
        \]
        \vspace{0.4cm}
        
        \textbf{Variable Definitions:}
        \begin{itemize}
            \item $x$ \textbf{(Input):} Input feature vector or matrix.
            \item $\theta$ \textbf{(Parameters):} Learnable weights and biases.
            \item $y$ \textbf{(Output):} Predicted output or representation.
        \end{itemize}
    }

    \block{3. Step-by-Step Workflow}{
        Explain the step-by-step algorithm or procedure flow.
        
        \vspace{0.5cm}
        \textbf{Step 1: Input Processing} \\
        Describe how input data is tokenized, normalized, or embedded.
        
        \vspace{0.5cm}
        \textbf{Step 2: Core Operation} \\
        Explain the transformations, matrix operations, or neural layers.
        
        \vspace{0.5cm}
        \textbf{Step 3: Output Generation} \\
        Explain how the final score, probability, or context is computed.
    }

    % ================= COLUMN 2 =================
    \column{0.33}

    \block{4. Diagrams \& Architecture}{
        \begin{center}
            % Replace the placeholder below with your actual diagram:
            % \includegraphics[width=0.7\linewidth]{architecture.png}
            \vspace{1cm}
            \fbox{\parbox[c][6cm][c]{0.7\linewidth}{\centering \Large \textbf{[ Insert Your Diagram / Figure Here ]}}}
            \vspace{1cm}\\
            \textit{Figure 1: Architectural diagram and computational pipeline.}
        \end{center}
    }

    \block{5. Complexity / Performance Analysis}{
        Detail the computational complexity, efficiency, and resource bottlenecks.

        \vspace{0.4cm}
        \textbf{1. Time Complexity Breakdown: $\mathcal{O}(\dots)$}
        \begin{itemize}
            \item \textbf{Operation A:} $\mathcal{O}(N \cdot d)$
            \item \textbf{Operation B:} $\mathcal{O}(N^2)$ (Primary bottleneck)
        \end{itemize}
        
        \vspace{0.4cm}
        \textbf{2. Space (Memory) Complexity: $\mathcal{O}(\dots)$}
        \begin{itemize}
            \item Intermediate storage requirements and cache behavior.
        \end{itemize}
    }

    \block{6. Real-World Applications}{
        Explain the practical use cases where this algorithm or model excels.
        
        \begin{itemize}
            \item \textbf{Application 1:} Use case description and benefits.
            \item \textbf{Application 2:} Edge cases and practical impact.
        \end{itemize}
    }

    % ================= COLUMN 3 =================
    \column{0.34}

    \block{7. Case Study / Simulation}{
        Present a concrete toy example, simulation, or walkthrough.
        
        \vspace{0.3cm}
        \textbf{Input Scenario:} Describe sample input values.
        
        \vspace{0.3cm}
        \textbf{Observed Outcome:} Describe what happens and why it validates your approach.
    }

    \block{8. Comparative Analysis}{
        \normalsize \color{black}
        \vspace{0.5cm}
        \begin{center}
        \resizebox{\linewidth}{!}{
        \renewcommand{\arraystretch}{1.8}
        \begin{tabular}{l l l l}
        \toprule
        \textbf{Feature / Property} & \textbf{Traditional Approach} & \textbf{Alternative Approach} & \textbf{Proposed Method} \\
        \midrule
        Processing Style & Sequential & Recurrent & Parallelized \\
        Complexity & $\mathcal{O}(N)$ & $\mathcal{O}(N \cdot d)$ & $\mathcal{O}(N^2 \cdot d)$ \\
        Scalability & Low & Moderate & Very High \\
        Performance & Baseline & Good & \textbf{State of the Art} \\
        \bottomrule
        \end{tabular}
        }
        \end{center}
    }

    \block{9. References}{
        \large \color{black}
        \begin{itemize}
            \item [1] Author, A., et al. (Year). Paper Title. \textit{Conference / Journal Name}.
            \item [2] Author, B., et al. (Year). Paper Title. \textit{Conference / Journal Name}.
        \end{itemize}
    }

\end{columns}

\end{document}
